\section{Economic interpretation and the next decision problem}
\label{sec:discussion125}

The economically relevant output is a feasible investment and its value
under the stipulated recursive preference. A neural continuation can be
useful in this account in two ways: it can propose an action that reduces
original-policy loss, or it can reduce the complete information cost needed
to reach a specified economic tolerance. A reduction in residual span,
critic training error or an optimizer's internal objective establishes
neither consequence without the accompanying action and cost evidence.

Table~\ref{tab:economics125} puts the improvement in original economic
units. Subtracting its outward-rounded intervals, the warm adaptive
service reduces mean cost relative to the completed fused control by
between $0.0002419$ and $0.0002523$ on this census. Its mean cost minus
that of the direct original online control lies between $-0.0000007$
and $0.0000028$. The comparison therefore does not certify an economic
improvement over that strong conventional control. These are objective
units in the specified model, not an estimated monetary benefit.
The eight direct-online cases consume $0.3750481$ complete primary
seconds in total, or about $0.0468810$ seconds per decision, with no
training acquisition bill and all eight stronger targets attained.
They are a subset of the fused-control row, not additional observations
or additional work. This one observed clock is descriptive; the full
decision cost and adoption conditions remain part of the comparison.

Complete Bellman supports make this distinction mathematical. A paid
calculation asserts a restriction on all feasible actions. Its information
should be reused in its full form when comparing decisions. The exact
finite representation shows when a numerical certificate is the strongest
one compatible with that declared record, and when apparently adequate
nodewise checks admit an impossible global support. Its rational action
example strengthens a guaranteed gain by $3/8$ in the specified
information class. A conventional direct-support calculation also obtains
that example's gain. The theorem establishes the full compatibility
problem; practical superiority of a complete-support solver over strong
conventional bounds requires a separate executed comparison.

The precision account treats new information as an economic input. A
smaller optimizer bracket may warrant a tighter risk calculation, whereas
a fixed terminal precision can leave a service unable to resolve the
action. Width, regret and improvement have different gates: a narrow
derivative bracket alone does not provide an absolute value enclosure,
and an optimal-future improvement must still pay for the actual future
controller. The implemented finite ladder records these distinctions.
It preserves qualified decisions when a new calculation cannot certify
a better one; unsuccessful acquisition remains part of the cost.


A frozen service outcome exposes this distinction directly. One warm
and fixed-precision cell reaches the capacity action $A=3/16$ and a
singleton original optimizer bracket, but its stored value-derived regret
bound is approximately $1.780734\times10^{-5}$, above the declared
$1/65536$ target. A separately recorded post-freeze rational diagnostic
reconstructs the original boundary derivative from the saved supports and
proves $q'(A)<0$. Strong convexity therefore identifies $A$ as the unique
constrained optimizer. The unchanged future controller then gives regret
at most $\delta=15/2^{24}$, approximately $8.94\times10^{-7}$, by
Theorem~\ref{thm:paidwork125}. The action and its absolute value interval
do not change. This diagnostic isolates an unnecessarily weak decision
gate and supplies a precise prospective implementation correction; it does
not alter the frozen target counts, assert a new policy gain, or claim
that the improved gate was deployed in the experiment.

The representation comparison is equally explicit. An affine--ReLU circuit
and the identical native hinge function are alternative encodings of the
same acquired decision rule. Their shared geometry does not identify a
neural-specific statistical or computational effect. A convincing neural
advantage must survive comparisons with conventional representations
having the same acquired labels, features, feasible policies and work
budget, and with suitable analytic or semianalytic original-model solvers.
The retained quadratic wins and empty adoption regions are evidence about
that requirement. They are not removed when a new certificate improves
an internal bound.

The current application is a stylized constrained investment economy.
Its continuous innovation and nested tail preference make qualification
nontrivial, but it has one control, one innovation factor and two decision
dates. The derivative constants, exact interpolation and execution results
have this stated scope. Neither a black-box continuous-risk lower bound
nor a scalar support-information minimax result excludes a faster method
that uses additional primitive formulas. Nor does a conditional precision
cost law prove a uniform lower bound over all economic algorithms.

The research programme therefore has concrete remaining tests. A neural
method must acquire action-relevant continuation information more cheaply
than equally informed strong conventional methods, including an independent
refit account when its training rule changes. A multidimensional extension
must replace scalar ordering and control the cost of continuous risk,
constraints and boundary interactions. Economic adoption needs either a
calibrated decision environment or an explicit range of values of the
computational resource. A new actor's population claim requires its own
independent sample and inference, rather than reuse of an earlier actor's
confidence interval. These are extensions of the original NBO objective
and research ambition, with the present theorem, code and preserved
negative evidence specifying what a successful extension must improve.
